% Vantagens_participar_SPE.tex
\documentclass[11pt,a4paper]{article}
\usepackage[utf8]{inputenc}
\usepackage[T1]{fontenc}
\usepackage[portuguese]{babel}
\usepackage{lmodern}
\usepackage{geometry}
\geometry{margin=2.5cm,headheight=15pt}
\usepackage{microtype}
\usepackage{xcolor}
\usepackage{titlesec}
\usepackage{fancyhdr}
\usepackage{hyperref}
\usepackage{enumitem}
\usepackage{booktabs}
\usepackage{verbatim}
\usepackage{setspace}
\usepackage{tcolorbox}
\usepackage{tikz}
\onehalfspacing

% Define professional color palette
\definecolor{spe-teal}{RGB}{0,102,102}
\definecolor{spe-accent}{RGB}{255,102,0}
\definecolor{spe-light}{RGB}{230,245,245}
\definecolor{spe-dark}{RGB}{20,40,60}

% Enhanced section formatting
\titleformat{\section}{\bfseries\Large\color{white}\filcenter}{}{0pt}{%
  \colorbox{spe-teal}{\parbox{\dimexpr\textwidth-2\fboxsep\relax}{\thesection\ #1}}}
\titlespacing{\section}{0pt}{12pt}{8pt}

\titleformat{\subsection}{\bfseries\large\color{spe-teal}}{\thesubsection}{0.8em}{}
\titlespacing{\subsection}{0pt}{10pt}{6pt}

\hypersetup{colorlinks=true,linkcolor=spe-accent,urlcolor=spe-accent}

% Fancy header with visual separator
\pagestyle{fancy}
\fancyhf{}
\fancyhead[L]{\small\color{spe-dark}\textbf{SPE ISPTEC Student Chapter}}
\fancyhead[R]{\small\color{spe-dark}Guia de Vantagens}
\renewcommand{\headrulewidth}{2pt}
\renewcommand{\headrule}{{\color{spe-accent}\hrule\@height\headrulewidth\@width\headwidth}}
\fancyfoot[C]{\small\color{spe-dark}\thepage}

% List formatting
\setlist[itemize]{leftmargin=1.5em,itemsep=4pt,parsep=2pt}
\setlist[enumerate]{leftmargin=1.5em,itemsep=4pt,parsep=2pt}

\title{}
\author{}
\date{}

\begin{document}

% Custom title page
\vspace*{1cm}
\begin{center}
{\color{spe-teal}\rule{\textwidth}{3pt}}\\
\vspace{0.5cm}
{\Huge\bfseries\color{spe-dark} Vantagens de Participar}\\
{\huge\bfseries\color{spe-accent} na SPE e no Student Chapter}\\
\vspace{0.3cm}
{\LARGE\itshape\color{spe-teal} Guia Prático para Estudantes}\\
\vspace{0.5cm}Participar — a Mensagem}
\textcolor{spe-dark}{Participar na SPE e no Student Chapter do ISPTEC transforma o teu percurso académico em vantagens concretas.} Ao juntares-te, tens acesso a conhecimento técnico relevante, contactos profissionais que facilitam estágios e empregos, oportunidades para apresentar trabalhos e experiências de liderança que enriquecem o teu currículo. 

\textbf{\color{spe-accent}É um investimento de tempo de baixo risco e alto retorno:} a participação regular aumenta a empregabilidade, a visibilidade académica e o acesso a bolsas.

\begin{tcolorbox}[colback=white,colframe=spe-accent,arc=3mm,boxrule=2pt,left=8pt,right=8pt,top=8pt,bottom=8pt]
{\bfseries\color{spe-dark}Benefícios-chave:}
\begin{itemize}
  \item[\color{spe-accent}✓] Acelera a entrada no mercado de trabalho
  \item[\color{spe-accent}✓] Fornece competências práticas aplicáveis imediatamente
  \item[\color{spe-accent}✓] Proporciona experiências de liderança e organização
  \item[\color{spe-accent}✓] Eleva a tua visibilidade académica e profissional
  \item[\color{spe-accent}✓] Abre portas a bolsas, visitas técnicas e oportunidades exclusivas
\end{itemize}
\end{tcolorbox}

\vspace{1em}
\begin{center}
\fbox{\textbf{\color{spe-accent}Regista-te agora:} \url{https://www.spe.org/join/}}
\end{center}

\section{Sumário executivo}
Participar na SPE dá acesso a uma comunidade profissional, recursos técnicos e oportunidades práticas que aceleram a carreira académica e profissional. Um Student Chapter acrescenta apoio local, financiamento e visibilidade.

\section{Porquê participar — mensagem persuasiva}
Participar na SPE e no Student Chapter do ISPTEC transforma o teu percurso académico em vantagens concretas. Ao juntares--te, tens acesso a conhecimento técnico relevante, contactos profissionais que facilitam estágios e empregos, oportunidades para apresentar trabalhos e experiências de liderança que enriquecem o teu currículo. Trata--se de um investimento de tempo de baixo risco e alto retorno: a participação regular aumenta a empregabilidade, a visibilidade académica e o acesso a bolsas e oportunidades.

\begin{itemize}
  \item Acelera a entrada no mercado de trabalho.
  \item Fornece competências práticas aplicáveis imediatamente.
  \item Proporciona experiências de liderança e organização.
  \item Eleva a tua Gerais da SPE}

\begin{enumerate}[label=\color{spe-accent}\textbf{\arabic*.}]
  \item[\color{spe-accent}\textbf{🤝}] \textbf{Networking com a Indústria} — Conecta com profissionais, ex-alunos e potenciais empregadores
  \item[\color{spe-accent}\textbf{📚}] \textbf{Acesso a Literatura Técnica} — Base de dados OnePetro com milhares de artigos
  \item[\color{spe-accent}\textbf{💼}] \textbf{Formação Profissional} — Webinars, cursos e certificações especializadas
  \item[\color{spe-accent}\textbf{🏆}] \textbf{Competições Académicas} — PetroBowl, PetroGames e desafios internacionais
  \item[\color{spe-accent}\textbf{💰}] \textbf{Financiamento e Bolsas} — Grants para eventos, atividades e desenvolvimento
  \item[\color{spe-accent}\textbf{⭐}] \textbf{Credibilidade Profissional} — Reconhecimento global da SPE
  \item[\color{spe-accent}\textbf{🎁}] \textbf{Descontos Exclusivos} — Em revistas, eventos e serviços
  \item[\color{spe-accent}\textbf{🌍}] \textbf{Comunidade Global} — Conexão com estudantes em 140+ países
\end{enumeratxtbf{Acesso a literatura técnica.}
  \item \textbf{Formdo Student Chapter ISPTEC}

\begin{tcolorbox}[colback=spe-light,colframe=spe-teal,arc=3mm,boxrule=2pt]
\begin{itemize}[label=\color{spe-teal}\textbf{◆}]
  \item \textbf{Apoio Financeiro} para palestras, workshops e visitas técnicas
  \item \textbf{Suporte Organizacional} da SPE e da instituição
  \item \textbf{Acesso Prioritário} a competições nacionais e internacionais
  \item \textbf{Mentoria com Profissionais Locais} — orientação de carreira
  \item \textbf{Desenvolvimento de Liderança} — oportunidades de organizar eventos
  \item \textbf{Comunidade Próxima} — amigos, colegas e networking local
\end{itemize}
\end{tcolorbox
\section{Benefícios específicos do Student Chapter}
\begin{itemize}
  \item \textbf{Apoio financeiro para atividades locais.}
  \item \textbf{Suporte organizacional.}
  \item \textbf{Acesso prioritário a competições e eventos estudantis.}
  \item \textbf{Mentoria e contacto com profissionais locais.}
  \item \textbf{Desenvolvimento de liderança.}
\end{itemize}

\section{Como aproveitar estes benefícios — passo a passo para estudantes}
\begin{itemize}
  \item \textbf{Inscrever-te na SPE:} visita \url{https://www.spe.org/join/} e cria conta com o teu email institucional; escolhe a categoria \texttt{Student}.
  \item \textbf{Participar em eventos:} inscreve-te nas palestras e workshops do capítulo; guarda os certificados no teu portfólio.
  \item \textbf{Formar ou integrar equipa (PetroBowl/PetroGames):} reúne colegas e treina regularmente.
  \item \textbf{Candidatar-te a bolsas e grants:} solicita apoio do Chapter Advisor e prepara um orçamento simples.
  \item \textbf{Usar OnePetro e os recursos SPE:} pesquisa artigos relevantes para os teus trabalhos.
  \item \textbf{Voluntariar-te em eventos:} oferece 2–3 horas para logística para ganhar experiência e networking.
  \item \textbf{Guardar certificados e construir portfólio:} cria pasta \texttt{SPE-Certificados} no Google Drive com os PDFs de participação.
  \item \textbf{Pedir mentoria:} contacta ex-alunos ou profissionais via email/LinkedIn solicitando alguns minutos de orientação.
\end{itemize}

\section{Mensagens e pitches rápidos}
\begin{itemize}
  \item UCronograma para o 1.º Ano — Foco em Alto Impacto}

Estratégia: atividades de baixo custo, máxima qualidade e participação.

\begin{tcolorbox}[colback=white,colframe=spe-accent,arc=3mm,boxrule=2pt]
\textbf{\color{spe-dark}Trimestre 1:} Construir base e confiança
\begin{itemize}[label=\color{spe-accent}→]
  \item 2 palestras com profissionais (networking + aprendizado)
  \item 1 workshop de soft skills (CV, entrevistas)
  \item 1 atividade social de integração (conhecer colegas)
\end{itemize}— Próximos Passos}

\begin{tcolorbox}[colback=spe-light,colframe=spe-teal,arc=3mm,boxrule=2pt]
\begin{itemize}[label=\color{spe-teal}☑]
  \item[\color{spe-teal}☐] Criar ficheiro \texttt{membros.csv} (nome, email, curso, ano, membroSPE, data\_inscricao)
  \item[\color{spe-teal}☐] Confirmar conta SPE na plataforma
  \item[\color{spe-teal}☐] Guardar certificados e documentos no Google Drive
  \item[\color{spe-teal}☐] Registar-se nos próximos eventos e voluntariar-se
  \item[\color{spe-teal}☐] Contactar Chapter Advisor para candidaturas a grants
  \item[\color{spe-teal}☐] Partilhar este guia com colegas interessados
\end{itemize}
\end{tcolorboxor{spe-dark}Trimestre 2:} Dinamizar e competir
\begin{itemize}[label=\color{spe-accent}→]
  \item Mesa-redonda com ex-alunos (trajetórias de carreira)
  \item PetroBowl interno ou treino de equipa
  \item Sessão técnica com profissional
\end{ite2em}

\begin{center}
{\color{spe-teal}\rule{0.6\textwidth}{2pt}}\\
\vspace{1em}
{\large\bfseries\color{spe-dark}Junte-se à comunidade SPE ISPTEC}\\
\vspace{0.5em}
{\color{spe-accent}\textbf{Transforme a sua carreira — registe-se em https://www.spe.org/join/}}\\
\vspace{1em}
{\color{spe-teal}\rule{0.6\textwidth}{2pt}}
\end{center}

\vspace{0.5em}

\begin{tcolorbox}[colback=white,colframe=spe-accent,arc=3mm,boxrule=2pt]
\textbf{\color{spe-dark}Trimestre 3:} Consolidar e expandir
\begin{itemize}[label=\color{spe-accent}→]
  \item Repetir atividades de maior adesão
  \item 1 visita técnica (refinaria, laboratório ou empresa local)
  \item Planear atividades para o ano seguinte
\end{itemize}
\end{tcolorbox}

\textbf{\color{spe-dark}Implementação Prática:}
\begin{itemize}
  \item Usar Google Forms para inscrições (automático)
  \item Templates SPE para agendas (consistência visual)
  \item Gerar certificados PDF automáticos em Google Sheets
\end{itemize}orkshop CV com RH local.)}
  \item \textbf{Trimestre 2:} Mesa-redonda + PetroBowl interno + sessão com ex-alunos. \textit{(Ex.: organizar equipa representante e treino quinzenal.)}
  \item \textbf{Trimestre 3:} Repetir atividades de maior adesão; preparar 1 visita técnica local opcional (avaliar orçamento). \textit{(Ex.: visita à refinaria ou laboratório universitário.)}
\end{enumerate}

Implementação rápida: usar formulários Google para inscrições; templates SPE para agendas; certificados PDF automáticos (gerador simples em Google Forms/Sheets).

\section{Checklist rápido}
\begin{itemize}
  \item Criar ficheiro \texttt{membros.csv} com colunas: nome, email, curso, ano, membroSPE, data\_inscricao, renovado.
  \item Confirmar conta SPE e guardar certificados no Google Drive.
  \item Registar-se nos próximos eventos do capítulo e, se possível, voluntariar-se.
  \item Pedir apoio ao Chapter Advisor para candidaturas a bolsas ou grants.
\end{itemize}

\section{Fontes e leituras úteis}
\begin{itemize}
  \item SPE — Homepage e secção de membros: \url{https://www.spe.org/}
  \item OnePetro (biblioteca técnica): \url{https://onepetro.org/}
  \item Consultar a secção "Student Chapters" no site da SPE para guias específicos e Student Chapter Grants.
\end{itemize}

\vspace{1em}
\noindent\textbf{Próximos passos sugeridos:} preparar o pacote de comunicação (flyer + email + slide), agendar a primeira palestra e lançar a campanha de inscrições.

\end{document}
